\begin{center}
\begin{tikzpicture}[x=.88cm,y=.88cm,font=\small]
  \fill[paper] (0,0) rectangle (12.3,6.65);

  \draw[source!70,line width=.8pt] (.5,3.04) .. controls (3.4,3.42) and
    (7.0,2.5) .. (11.85,3.08);
  \node[text=source,below,font=\scriptsize] at (8.25,2.86) {黄河方向（示意）};
  \draw[source!55,dashed] (.62,1.34) .. controls (3.7,.86) and
    (7.3,1.9) .. (11.72,1.18);
  \node[text=source,below,font=\scriptsize] at (8.2,1.09) {长江方向（示意）};

  \fill[dynasty!44] (6.15,4.66) ellipse (1.62 and .96);
  \node[align=center,text=white] at (6.15,4.66) {北魏\\平城—洛阳中枢};

  \fill[timeline!27] (2.05,4.75) ellipse (1.34 and .8);
  \node[align=center] at (2.05,4.75) {河西、柔然\\西北与北部边境};

  \fill[source!27] (9.98,4.82) ellipse (1.34 and .82);
  \node[align=center] at (9.98,4.82) {山东、河北\\州郡与军镇};

  \fill[dynasty!34] (7.98,1.98) ellipse (1.55 and .82);
  \node[align=center,text=white] at (7.98,1.98) {南朝宋—梁\\建康朝廷};

  \fill[timeline!25] (3.52,1.88) ellipse (1.44 and .8);
  \node[align=center] at (3.52,1.88) {荆、湘、益\\长江上游军镇};

  \draw[->,very thick,color=dynasty]
    (6.78,3.84) .. controls (7.38,3.24) and (7.7,2.86) .. (7.85,2.67);
  \node[text=dynasty,right,font=\scriptsize,align=left] at (7.27,3.31)
    {淮、汉前线与\\南北交涉};

  \draw[->,thick,dashed,color=timeline]
    (5.14,4.48) .. controls (4.44,4.43) and (3.85,4.56) .. (3.42,4.68);
  \node[text=timeline,above,font=\scriptsize,align=center] at (4.22,5.1)
    {军镇、边贸与\\北部防务};

  \draw[->,thick,dashed,color=source]
    (4.68,2.1) .. controls (5.52,2.2) and (6.08,2.18) .. (6.42,2.12);
  \node[text=source,above,font=\scriptsize,align=center] at (5.55,2.5)
    {上游军镇、粮道\\与建康的联系};

  \draw[->,thick,color=timeline]
    (7.1,4.88) .. controls (7.82,4.87) and (8.24,4.85) .. (8.56,4.84);
  \node[text=timeline,above,font=\scriptsize] at (7.86,5.22)
    {州郡、户籍与交通};

  \node[draw=source!75,fill=paper,rounded corners=2pt,align=center,
    inner sep=3pt,font=\scriptsize] at (10.05,.6)
    {南北“对峙”并非一条固定边界：都城、军镇、河道与州郡网络\\
    在不同阶段以不同方式连接或隔开两侧};
\end{tikzpicture}

\smallskip
\small\textit{图\quad 北魏与南朝之间的几个政治、军政和交通重心，约5世纪至6世纪初：据《魏书·世祖纪》《高祖纪》《世宗纪》、沈约《宋书》、姚思廉《梁书》及《资治通鉴》卷一百二十六至一百四十八，概括平城、洛阳、建康、荆湘与黄河、长江之间的关系。政权和地点跨越长时段列示，不表示同一时点的完整形势；色块不绘疆界，箭线亦非战线或行政控制范围。}
\end{center}

\phantomsection\label{fig:northern-southern-courts}
